% The paper's central diagram. Hand-authored TikZ rather than an image so it sets in the
% document's own fonts and survives a style change. Both panels are deliberately identical
% except for the loss and the direction of the second arrow, because that is the whole point.
%
% Geometry note: the vertical gap between the two boxes has to be wide enough for the two
% arcs to read as a loop. A gap smaller than roughly the box height turns them into stubs.
\begin{figure}[t]
\centering
\begin{tikzpicture}[
  box/.style  ={draw=black!22, fill=black!4, rounded corners=4pt,
                minimum width=2.6cm, minimum height=0.88cm, font=\normalsize},
  hdr/.style  ={font=\small\bfseries},
  eqn/.style  ={font=\small},
  side/.style ={font=\scriptsize, text=black!55},
  vmark/.style ={font=\small\bfseries},
  fwd/.style  ={-{Latex[length=2.2mm]}, line width=0.9pt, black!78},
  dead/.style ={-{Latex[length=2.2mm]}, line width=0.9pt, black!30, densely dashed},
  live/.style ={-{Latex[length=2.2mm]}, line width=1.2pt, flipteal},
]

% ---------------- panel (a): the published intervention ----------------
\begin{scope}[shift={(0,0)}]
  \node[box] (ax)  at (0,2.05) {$x$};
  \node[box] (axo) at (0,0)    {$\tilde x$};
  \node[hdr] at (0,3.40) {(a)\; add a contrastive objective};
  \node[eqn] at (0,2.85)
    {$\mathcal{L}=\mathcal{L}_{\mathrm{gen}}+\mathcal{L}_{\mathrm{pos}}
                 +\mathcal{L}_{\mathrm{neg}}$};
  \draw[fwd]  ([xshift=-7mm]ax.south)  to[bend right=45] ([xshift=-7mm]axo.north);
  \draw[dead] ([xshift=7mm]axo.north) to[bend right=45] ([xshift=7mm]ax.south);
  \node[side, anchor=east] at (-1.55,1.02) {obfuscate};
  \node[side, anchor=west, text=black!40] (adl) at (1.55,1.02) {deobfuscate};
  \node[vmark, text=failred, anchor=west] at (adl.east) {\,\texttimes};
\end{scope}

% ---------------- panel (b): reading the same pairs backwards ----------------
\begin{scope}[shift={(7.6,0)}]
  \node[box] (bx)  at (0,2.05) {$x$};
  \node[box] (bxo) at (0,0)    {$\tilde x$};
  \node[hdr, text=flipteal] at (0,3.40) {(b)\; flip the training pairs};
  \node[eqn] at (0,2.85) {$\mathcal{L}=\mathcal{L}_{\mathrm{gen}}$};
  \draw[fwd]  ([xshift=-7mm]bx.south)  to[bend right=45] ([xshift=-7mm]bxo.north);
  \draw[live] ([xshift=7mm]bxo.north) to[bend right=45] ([xshift=7mm]bx.south);
  \node[side, anchor=east] at (-1.55,1.02) {obfuscate};
  \node[side, anchor=west] (bdl) at (1.55,1.02) {deobfuscate};
  \node[vmark, text=flipteal, anchor=west] at (bdl.east) {\,\checkmark};
\end{scope}

\end{tikzpicture}
\caption{The attribution this paper tests. Both panels fine-tune the same model on the same
obfuscation pairs. \textbf{(a)}~Contrastive Fine-Tuning augments the generation loss with two
equivalence-judgement terms and is credited with restoring the reverse direction. In our
experiments it does not. Reverse success stays at zero and falls below the untouched model.
\textbf{(b)}~Keeping the plain generation loss and reading half of the same pairs backwards
restores the reverse direction at identical training cost. The recovery comes from the
direction of the data, not from the objective.}
\label{fig:attribution}
\end{figure}
